% Part: many-valued-logic
% Chapter: three-valued-logics
% Section: introduction

\documentclass[../../../include/open-logic-section]{subfiles}

\begin{document}

\olfileid{mvl}{thr}{int}

\olsection{جان پہچان}

جے اسیں اک ہور قدر~$\Undef$، $\True$ تے $\False$ نال
شامل کریے، تاں تِن قدراں والی منطق ملدی اے۔ بھاویں صداقتی
قدر اکّو ودھدی اے، فیر وی $\lnot$، $\land$، $\lor$ تے $\lif$
دے صداقتی فنکشن متعین کرن دے کئی طریقے نیں۔ منطق وچ ایہناں
فنکشوناں دی کوئی وی ترکیب ورتی جا سکدی اے۔ نالے ایہ چُن سکدے
آں کہ صرف $\True$ مقرر ہووے، یا $\True$ تے~$\Undef$ دوویں۔

ایتھے اسیں تِن قدراں والیاں سبھ توں معروف منطقاں وچوں کجھ
منطقاں، اوہناں دے بناؤن دے سبب تے کجھ خاصیتاں پیش کردے آں۔

\end{document}
